% Generated draft: requires mathematical and source review.
\documentclass[11pt]{article}
\usepackage[margin=25mm]{geometry}
\usepackage{fontspec}
\setmainfont{Latin Modern Roman}
\usepackage{amsmath,amssymb,mathtools,unicode-math}
\setmathfont{Latin Modern Math}
\usepackage{enumitem,needspace,xurl,hyperref}
\setlength{\parindent}{0pt}
\setlength{\parskip}{5pt}
\setlength{\emergencystretch}{3em}
\newcommand{\N}{\mathbb{N}}
\newcommand{\Z}{\mathbb{Z}}
\newcommand{\Q}{\mathbb{Q}}
\newcommand{\R}{\mathbb{R}}
\newcommand{\C}{\mathbb{C}}
\newcommand{\F}{\mathbb{F}}
\newcommand{\A}{\mathbb{A}}
\newcommand{\PP}{\mathbb{P}}
\newcommand{\E}{\mathbb{E}}
\newcommand{\eps}{\varepsilon}
\newcommand{\dd}{\,\mathrm{d}}
\begin{document}
\section*{UnsolvedMath ID: 30005651}
\textbf{Source problem code:} catalogue entry
\subsection{Definitions and mathematical statement}
The supplied question concerns a finite group \(G\) admitting an irreducible representation \(V\) with \(\dim V>1\). The field of definition of \(V\), the base scheme for the smooth \(G\)-schemes, and any restrictions on their characteristics are not specified. These data must be fixed before the question has a precise domain.

A conditional formulation of the intended stable-connectivity question is as follows. Fix a base scheme \(S\), an intended category \(\mathrm{Sm}^{G}_{S}\) of smooth \(G\)-schemes over \(S\), and the intended Nisnevich topology \(\tau\). For an ordinary spectrum-valued presheaf \(E\) on this category, write
\[
\underline{\pi}^{\tau}_{n}(E)=a_{\tau}\bigl(X\longmapsto\pi_n(E(X))\bigr),
\]
where \(a_{\tau}\) denotes sheafification. Under the interpretation that ``Nisnevich-locally connected'' means connective in the stable sense, this condition is
\[
\underline{\pi}^{\tau}_{n}(E)=0\qquad\text{for every integer }n<0.
\]
Let \(L^{G}_{\mathbb A^{1}}\) denote the intended equivariant \(\mathbb A^{1}\)-localization. With these conventions, the candidate mathematical question is whether, for every such finite group \(G\) and every such presheaf \(E\),
\[
\left[\underline{\pi}^{\tau}_{n}(E)=0\ \text{for all }n<0\right]
\quad\Longrightarrow\quad
\left[\underline{\pi}^{\tau}_{n}\!\left(L^{G}_{\mathbb A^{1}}E\right)=0\ \text{for all }n<0\right].
\]

This is a conditional reconstruction, not a uniquely determined formal statement from the supplied evidence. The record does not specify whether the spectra are ordinary, genuine equivariant, or motivic spectra; the relevant stabilization and homotopy sheaves may therefore differ from those displayed above. It also does not specify the equivariant Nisnevich topology, the action on the affine-line interval, or whether the localization includes Nisnevich localization. The word ``connected'' is not defined: the negative-degree vanishing condition above expresses stable connectiveness, whereas other uses of connectedness can impose a different condition in degree zero. The original supplied excerpt asks only whether stable connectivity still holds and does not resolve these choices.
\subsection{Short English statement}
For a finite group admitting an irreducible representation of dimension greater than one, the supplied question asks whether equivariant \(\mathbb A^{1}\)-localization preserves Nisnevich-local stable connectivity of spectral presheaves on smooth \(G\)-schemes. A conditional formulation uses vanishing of negative homotopy sheaves, but the base, representation field, equivariant spectral category, topology, and connectivity convention require clarification. The supplied registry labels the problem open; this status has not been independently verified.
\subsection{Sources}
\begin{itemize}
\item Local UnsolvedMath registry, permanent ID \(30005651\), title: ``Equivariant Stable Connectivity in Higher Representation Dimension,'' source problem code \texttt{catalogue entry}; collection: Oberwolfach Reports Open Problems. No individual proposer is supplied. The proposed year is recorded as 2023.
\item Supplied report citation: \emph{Tensor-Triangular Geometry and Interactions} (2024), Oberwolfach Reports, pp.\ 2303--2358, DOI \texttt{10.4171/owr/2023/40}. \url{https://doi.org/10.4171/owr/2023/40}. The supplied original excerpt asks whether stable connectivity still holds; no narrower passage location or full surrounding mathematical setup is supplied.
\item Supplied literature source: Bachmann, ``motivic spectral Mackey functors.'' \url{https://arxiv.org/abs/2205.13926}. No specific supporting passage or theorem is identified in the record.
\item Literature checked at: \texttt{2026-08-22}. This is supplied dated curation, not an independently verified assessment. The supplied assessment labels the problem open and states that the report presents stable connectivity as an open question beyond elementary examples and gives no general proof for groups with higher-dimensional irreducible representations. The background explicitly requests independent verification of time-sensitive status and cited claims.
\item The registry assigns the statement status \texttt{corrected\_verified} and reports that the official report was checked directly and that the clean statement restores definitions and preserves the mathematical target. These are supplied provenance claims; no source visit or independent verification is asserted here. The citation gives the report year as 2024, while the DOI, proposed year, and workshop-year tag refer to 2023; this distinction is retained without reconciliation.
\end{itemize}
\end{document}
